\section{Implementation-Facing Mathematical Properties}
\label{sec:theory}
\label{app:lean}

The statements below are formalized in \path{BAMOR/Implementation.lean},
namespace \path{BAMOR.Implementation}. They distinguish target construction,
feedback statistics, conditional tracking and archive evaluation. Algebra is
over exact real numbers. The Lean source and its imported auxiliary results
are replayed with Lean~4.28.0 and Mathlib~4.28.0.

\subsection{Soft targets and continuation masks}
\label{app:soft_target}

Let $u\in\Delta^{m-1}$, $b=\gamma(1-d_{\mathrm{term}})$ and
$e=-\alpha\log\pi_\theta(a'_{j^*}\mid s',v_{j^*},z')$.
For the selected vector $q$, define $Y_i=r_i+b(q_i+e)$.

\begin{proposition}[Consistent scalar projection]
\label{prop:soft_projection}
The v25 targets satisfy
\begin{equation}
 u^\top Y=u^\top r+b(u^\top q+e).
 \label{eq:soft_projection}
\end{equation}
At true termination $b=0$, $Y=r$. Replacing a continuing-task time-limit
mask $b=\gamma$ by zero removes $\gamma(q_i+e)$ from coordinate $i$.
\end{proposition}
\begin{proof}
Distribute the scalar product and use $\sum_i u_i=1$ for the common entropy
bonus. The mask identities follow by substitution. Lean:
\path{soft_target_projection}, \path{terminal_target}, and
\path{truncation_target_difference}.
\end{proof}

With the same selected candidate and entropy bonus in two targets,
$|q_i-\widetilde q_i|\leq\epsilon$ and $b\geq0$ give
$|Y_i-\widetilde Y_i|\leq b\epsilon$
(\path{frozen_target_lipschitz}). The qualification about selection matters.

\begin{proposition}[Candidate reselection can amplify vector differences]
\label{prop:reselection}
For $u=(1/2,1/2)$, consider candidates $(1,0),(0,0.9)$ and their perturbed
versions $(1,0),(0,1.1)$, with equal entropy bonuses. All candidate coordinates
change by at most $0.2$, but the first coordinate of the selected vector changes
by $1$. Both scalar maximizers are unique.
\end{proposition}
\begin{proof}
The scalar scores change from $(0.5,0.45)$ to $(0.5,0.55)$. The maximizing
candidate switches from the first to the second. The coordinatewise bound and
selected-vector difference are checked exactly by
\path{candidate_reselection_counterexample}; strict score inequalities are
checked by \path{candidate_reselection_unique}.
\end{proof}

Thus the scalar maximum's non-expansiveness cannot be transferred to its
selected vector. The counterexample also applies to equal twin critics and
fixed candidate actions; neither critic disagreement nor a tie is necessary.

\begin{proposition}[Margin-certified target stability]
\label{prop:candidate_margin}
Fix a nonempty finite set of candidate actions and their context inputs,
$u\in\Delta^{m-1}$, reward $r$ and bootstrap multiplier $b\geq0$.
For candidate $j$, let $q_j$ be the coordinatewise minimum of the two critic
vectors, $e_j$ its entropy bonus and $s_j=u^\top q_j+e_j$ its ranking score.
Compare these with perturbed values $\widetilde q_j,\widetilde e_j$ satisfying
\begin{equation}
 |q_{j,i}-\widetilde q_{j,i}|\leq\epsilon,\qquad
 |e_j-\widetilde e_j|\leq\rho,\qquad
 \delta=\epsilon+\rho,\quad \epsilon,\rho\geq0.
 \label{eq:candidate_perturbation}
\end{equation}
If $s_{j^*}-s_j>2\delta$ for every $j\neq j^*$, then $j^*$ is the unique
maximizer of both score sets. The corresponding soft targets satisfy
\begin{equation}
 |Y_i-\widetilde Y_i|\leq b\delta\quad\text{for every }i,
 \qquad
 |u^\top Y-u^\top\widetilde Y|\leq b\delta.
 \label{eq:margin_target_bound}
\end{equation}
\end{proposition}
\begin{proof}
The simplex conditions and the triangle inequality give
$|s_j-\widetilde s_j|\leq\epsilon+\rho=\delta$.
Hence $\widetilde s_{j^*}-\widetilde s_j
\geq s_{j^*}-s_j-2\delta>0$.
With the selected index unchanged, the coordinatewise target difference is
$b|(q_{j^*,i}-\widetilde q_{j^*,i})+(e_{j^*}-\widetilde e_{j^*})|
\leq b\delta$. Simplex projection gives the scalar bound.
Lean: \path{soft_score_perturbation},
\path{strict_margin_preserves_winner}, and
\path{margin_selected_target_lipschitz}.
\end{proof}

\subsection{Progress-score properties}
\label{app:progress}

For the first and second half-block reward means $a,b\in\mathbb R^m$ and
the frozen positive scale $s$, write
\begin{equation}
 P(a,b;s)=\frac{1}{m}\sum_i\tanh\frac{b_i-a_i}{s_i}.
\end{equation}

\begin{proposition}[Bounded, centered block feedback]
\label{prop:progress}
For $m\geq1$, $-1\leq P\leq1$. For any per-objective offset $c$,
$P(a+c,b+c;s)=P(a,b;s)$. Positive rescaling preserves the score if the
supplied scales are rescaled by the same coordinatewise factors. If
$q,P\in[-1,1]$ and $\eta\in[0,1]$, then
$(1-\eta)q+\eta P\in[-1,1]$.
\end{proposition}
\begin{proof}
Each hyperbolic tangent is in $(-1,1)$, so its finite average is in $[-1,1]$.
Offsets cancel in the difference, and common nonzero factors cancel in the
ratio. The score update is a convex combination. Lean:
\path{progress_score_bounds}, \path{progress_translation_invariant},
\path{progress_rescaling_invariant}, and \path{score_ema_bounds}.
\end{proof}

Starting at zero, v25's score EMA remains bounded by induction; unselected
arms and incomplete blocks leave scores unchanged. Its fixed $10^{-6}$ scale
floor means that rescaling invariance need not hold when clipping changes.
The propositions concern the statistic on supplied observations and scales,
not an invariant training trajectory. A reward increase caused by a mode
change enters the same numerator as an improvement caused by learning.

\subsection{Conditional tracking with update error}
\label{app:tracking_current}

Let $T_{n,o}$ be a family of maps on a metric space, $o_n$ the realized
selected arm and $q_n^*$ a reference trajectory. Assume, for all $n,o,x,y$,
\begin{align}
 d(T_{n,o}x,T_{n,o}y)&\leq\gamma d(x,y),\qquad 0\leq\gamma<1,\nonumber\\
 d(q_{n+1},T_{n,o_n}q_n)&\leq\epsilon,\nonumber\\
 d(T_{n,o_n}q_n^*,q_{n+1}^*)&\leq\Delta,
 \qquad \epsilon,\Delta\geq0.
 \label{eq:tracking_assumptions}
\end{align}
The last quantity includes the selected operator's mismatch with the reference,
not just physical environment drift.

\begin{proposition}[Approximate selected-map tracking]
\label{prop:tracking_current}
Under Eq.~\eqref{eq:tracking_assumptions}, for every $n\geq0$,
\begin{equation}
 d(q_n,q_n^*)\leq\gamma^n d(q_0,q_0^*)+
                    \frac{\epsilon+\Delta}{1-\gamma}.
\end{equation}
\end{proposition}
\begin{proof}
Two triangle inequalities yield
$d(q_{n+1},q_{n+1}^*)\leq\gamma d(q_n,q_n^*)+\epsilon+\Delta$.
Apply the geometric recurrence bound \path{BAMOR.anytime_tracking_bound}.
The full composition is \path{approximate_selected_tracking}.
\end{proof}

The residual term is an upper bound, not a lower error floor, and the initial
error term remains part of the bound. No comparison with CAPQL follows from
this recurrence alone.

The arm sequence may depend on observed history: the inequality compares two
inputs using the \emph{same} selected map at that step. It is not a comparison
between two feedback trajectories that can choose different arms. The finite
state/action auxiliary operator in \path{BAMOR.lean} uses per-objective action
maxima and frozen belief/penalty parameters. Its contraction theorem is valid
for that operator. Proposition~\ref{prop:reselection} shows why the vector
candidate update needs separate treatment.

\subsection{Continuous hypervolume and policy selection}
\label{app:archive_current}

Fix the mode, encoder, return vectors $J(p)$ and reference $r$. A candidate
$p$ denotes an actor/preference pair. For a finite archive $A$, define
\begin{equation}
 D_r(A)=\bigcup_{p\in A}[r,J(p)],\qquad
 \mathrm{HV}_r(A)=\operatorname{vol}_m(D_r(A)),
\end{equation}
where the interval is coordinatewise and is empty when any upper coordinate
is below the reference. This is the maximization convention.

\begin{proposition}[Continuous archive monotonicity]
\label{prop:hv_current}
The dominated region is compact and measurable with finite volume, and
$A\subseteq B$ implies
$D_r(A)\subseteq D_r(B)$ and $\mathrm{HV}_r(A)\leq\mathrm{HV}_r(B)$.
\end{proposition}
\begin{proof}
Each finite-coordinate rectangle is compact and measurable; the union is
finite, so its Lebesgue volume is finite. Every rectangle indexed by $A$ also
occurs in $B$, and Lebesgue measure is monotone. Lean:
\path{dominated_region_measurable}, \path{dominated_region_mono}, and
\path{hypervolume_mono}. The measure is initially valued in $[0,\infty]$;
\path{hypervolume_ne_top} proves finiteness, and
\path{hypervolume_toReal_mono} transfers monotonicity to the real-valued HV
used here. This concerns continuous volume rather than a grid-cell count.
\end{proof}

\begin{proposition}[From estimated archive scores to a selected policy]
\label{prop:selection_regret}
Fix a query and scalar scores $J,\widehat J$ on a nonempty finite archive $A$.
If $|\widehat J(p)-J(p)|\leq\epsilon$ for every $p\in A$ and
$\hat p\in\arg\max_{p\in A}\widehat J(p)$, then
\begin{equation}
 \max_{p\in A}J(p)-J(\hat p)\leq2\epsilon.
\end{equation}
\end{proposition}
\begin{proof}
For every $p\in A$,
$J(p)\leq\widehat J(p)+\epsilon\leq\widehat J(\hat p)+\epsilon
\leq J(\hat p)+2\epsilon$. Take the finite maximum. Lean:
\path{archive_selection_regret}.
\end{proof}

This is a conditional policy-selection bound, not regret of the training-arm
selector. It requires a uniform evaluation-error bound and an actual selection
rule. Archive HV supplies neither. The experiments report fixed-rollout archive
coverage and direct final-policy utility separately.

\subsection{Proof scope and retained auxiliary results}

The current results formalize target consistency, bounded score updates,
margin-certified target stability, conditional tracking and fixed-value archive
comparisons. The margin condition uses matched candidates and explicit bounds
on critic and entropy outputs; these bounds have not been established along
the training trajectories. Equality at the $2\delta$ threshold can produce a
tie, whereas failing the sufficient condition does not imply a switch.
The results do not prove
convergence of the implemented actor--critic or a causal improvement from its
progress selector. No uniform return-estimation error bound is established by
the three-reset evaluation protocol. Updating the encoder, re-estimating old
returns or replacing archive members changes the object in
Proposition~\ref{prop:hv_current}.

The earlier contraction, localized-drift, anchor and belief results remain in
\path{BAMOR.lean} and \path{RegimeBelief.lean}. Their definitions concern the
auxiliary frozen tabular model. Anchor concavity is with respect to the
preference, not actor parameters. \path{BAMOR-Counterproof.lean} retains its
Q-dependent-belief examples; its posterior-ratio threshold assumes a specific
power sequence, rather than deriving a BOCD detection rate. These results are
retained in the proof repository without duplicating their full historical
exposition in this appendix.

\subsection{Historical v23 target and score}
\label{app:v23_legacy}

The v23 critic used the episode-end mask $d$, including time limits, and
different scalar and vector targets:
\begin{equation}
 y_s=u^\top r+\gamma(1-d)(u^\top q+e),\qquad
 \mathbf y_v=r+\gamma(1-d)q.
\end{equation}
Hence $y_s-u^\top\mathbf y_v=\gamma(1-d)e$ for simplex weights. V24 repaired
both the mask and this target inconsistency; v25 retains those repairs.

For a 1,500-transition block, v23 used $g_t=w_t^\top r_t$ and
$R_k=\overline g_k-0.08\,\operatorname{std}(g_t)
-0.04(\overline g_k-\min_t g_t)$.
With $\delta_k=R_k-B$, its updates were $B\leftarrow0.9B+0.1R_k$,
$S\leftarrow0.9S+0.1|\delta_k|$, followed by
$q_o\leftarrow0.87q_o+0.13\tanh(\delta_k/\max(S,10^{-6}))$.
The first block initialized $B=R_k,S=1$ before these updates.
V25 replaces this score and starts its scored blocks after 5,000 transitions;
the operator configurations and family-allocation rule are unchanged.
